% ID: SEL-D0060-OA-E02
% ESTADO: cotejada. Segunda pasada visual y borrador compilado revisados.
% Fuente: Selectividad/sel_2011_matematicas_aplicadas/223A-EXAMEN 3.pdf
% SHA-256: 5998699e0b2992b988a8e0ed0d2e3581dd8fac52d6f9802587ebd40fee5ef608
% Página: 1; opción A; ejercicio original 2.

El beneficio, en miles de euros, alcanzado en una tienda de ropa el pasado año, viene dado por la función $B(t)$ expresada a continuación
\[
B(t)=\begin{cases}\dfrac{1}{8}t^2-t+5&\text{si }0\leq t\leq6\\[0.8em]\dfrac{t+1}{2}&\text{si }6<t\leq12\end{cases},
\]
$t$ es el tiempo transcurrido en meses.
\begin{enumerate}[label=\alph*)]
\item \textbf{(1 punto)} Estudie la derivabilidad de la función al cabo de 6 meses.
\item \textbf{(0.5 puntos)} \textquestiondown{}Cuándo fue mínimo el beneficio? \textquestiondown{}Cuál fue dicho beneficio?
\item \textbf{(1 punto)} Represente gráficamente la función $B(t)$. \textquestiondown{}Cuándo fue máximo el beneficio? \textquestiondown{}A cuánto ascendió?
\end{enumerate}
